\section{Introduction}
Myrtille est une charte graphique pensée pour produire des documents techniques avec un rendu clair et moderne. Son utilisation est faite pour être simple grâce à l'utilisation d'un préambule LaTeX standard. 

% ── SECTION 1 : COULEURS ───────────────────────────────────────
\section{Palette de Couleurs}

Cette section détaille l'ensemble des couleurs définies dans le préambule de Myrtille. Elles sont organisées par thématique pour faciliter leur utilisation.

\subsection{Couleurs Principales (Thème)}
Ces couleurs constituent l'identité visuelle de base et sont utilisées pour les éléments mis en évidence (titres, liens, bordures principales).

\vspace{10pt}
\colorswatch{accent}{0D6EFD}
\colorswatch{accentDark}{0A4FA8}
\colorswatch{accentPale}{E8F1FF}
\colorswatch{coverBg}{0F2A4A}

\subsection{Couleurs Neutres (Texte et Fonds)}
Ces nuances de gris et de couleurs pâles structurent le document, assurent la lisibilité du texte et définissent les arrière-plans secondaires.

\vspace{10pt}
\colorswatch{dark}{1A1D23}
\colorswatch{mid}{4A5060}
\colorswatch{light}{F0F2F5}
\colorswatch{border}{D0D6E0}
\colorswatch{sidebarBg}{F6F8FC}

\subsection{Statuts et Alertes}
Ces couleurs sont strictement réservées aux boîtes de messages et aux badges indiquant un état particulier (succès, avertissement, information).

\vspace{10pt}
\colorswatch{success}{198754}
\colorswatch{warn}{FFC107}
\colorswatch{warnDark}{664D03}
\colorswatch{warnBg}{FFF8E1}
\colorswatch{infoGreen}{0F5132}
\colorswatch{infoBg}{D1E7DD}

\subsection{Coloration Syntaxique (Code)}
Utilisées en arrière-plan et pour coloriser les éléments des blocs de code (mots-clés, chaînes de caractères, commentaires).

\vspace{10pt}
\colorswatch{codeBg}{F1F4F9}
\colorswatch{codeText}{1A2540}
\colorswatch{codeKw}{0550AE}
\colorswatch{codeStr}{1A7F37}
\colorswatch{codeCmt}{7A8599}
\colorswatch{inlineCode}{0D6EFD}


% ── SECTION 2 : TYPOGRAPHIE ET TITRES ──────────────────────────
\section{Typographie et Hiérarchie}

Le document utilise la police sans-serif (Helvetica) pour un rendu moderne et épuré.

\subsection{Ceci est un titre de Sous-section (\texttt{\textbackslash subsection})}
Le texte des paragraphes standards s'affiche dans la couleur noire par défaut du document. L'espacement entre les paragraphes est géré automatiquement sans indentation, offrant un confort de lecture optimal.

\subsubsection{Ceci est un titre de Sous-sous-section (\texttt{\textbackslash subsubsection})}
Les titres de niveau 3 utilisent la couleur \texttt{mid} pour marquer une séparation subtile dans le contenu sans surcharger visuellement la page. 

\vspace{10pt}
\hrulefull
\\La ligne ci-dessus est générée avec la commande \ci{\textbackslash hrulefull}, idéale pour séparer deux grandes idées au sein d'une même section.


% ── SECTION 3 : LISTES ET BADGES ───────────────────────────────
\section{Éléments de texte et Badges}

\subsection{Listes à puces et numérotées}

Voici un exemple de liste non ordonnée (\texttt{itemize}) avec différents niveaux :
\begin{itemize}
    \item Premier point de la liste principale.
    \item Deuxième point avec une sous-liste :
    \begin{itemize}
        \item Élément secondaire A.
        \item Élément secondaire B.
    \end{itemize}
    \item Troisième point.
\end{itemize}

Voici un exemple de liste ordonnée (\texttt{enumerate}) :
\begin{enumerate}
    \item Étape numéro un de la procédure.
    \item Étape numéro deux, qui suit logiquement la première.
    \item Étape finale du processus.
\end{enumerate}

\subsection{Badges et Code en ligne}

Les badges sont conçus pour qualifier rapidement des informations (type de variable, obligation, etc.).
Le code en ligne s'écrit avec la commande \ci{\textbackslash ci\{...\}}. Exemple : la variable \ci{config\_path} doit être définie.

\vspace{10pt}
\begin{tabularx}{\textwidth}{L{4cm} X}
    \toprule
    \rowcolor{tableHdr}
    \textbf{Commande} & \textbf{Rendu visuel et Usage} \\
    \midrule
    \ci{\textbackslash badgeReq} & \badgeReq ~ — Indique qu'un paramètre est obligatoire. \\
    \ci{\textbackslash badgeOpt} & \badgeOpt ~ — Indique qu'un paramètre est facultatif. \\
    \ci{\textbackslash badgeInt} & \badgeInt ~ — Type de donnée entier. \\
    \ci{\textbackslash badgeFloat} & \badgeFloat ~ — Type de donnée flottant. \\
    \ci{\textbackslash badgeStr} & \badgeStr ~ — Type de donnée chaîne de caractères. \\
    \ci{\textbackslash badgeMat} & \badgeMat ~ — Type matriciel ou array. \\
    \ci{\textbackslash badgePolar} & \badgePolar ~ — Coordonnées polaires. \\
    \bottomrule
\end{tabularx}


% ── SECTION 4 : CRÉATION DE BADGES PERSONNALISÉS ───────────────
\section{Création de Badges Personnalisés}

La charte graphique met à disposition une macro racine \ci{\textbackslash badge[couleur]\{texte\}} qui permet de générer de nouveaux badges en quelques secondes.

\subsection{À partir d'une couleur existante}
Si la couleur dont vous avez besoin figure déjà dans la palette (par exemple \texttt{infoGreen}), ajoutez simplement une nouvelle commande dans le préambule de votre document :

\begin{shblock}
\newcommand{\badgeNouveau}{\badge[infoGreen]{Nouveau}}
\end{shblock}

\textbf{Résultat visuel :} \badgeNouveau

\subsection{À partir d'une couleur inédite}
Bien que ce ne soit pas recommandé, si vous souhaitez utiliser une teinte qui n'est pas dans la charte standard, vous devez la définir avec \ci{\textbackslash definecolor}, puis créer le badge :

\begin{shblock}
% 1. Déclarer la couleur dans le préambule
\definecolor{customPurple}{HTML}{8E44AD}

% 2. Créer la commande associée
\newcommand{\badgePerso}{\badge[customPurple]{Spécial}}
\end{shblock}

\textbf{Résultat visuel :} \badgePerso


% ── SECTION 5 : BOÎTES D'INFORMATION ───────────────────────────
\section{Blocs d'Information et Notes}

Ces environnements permettent de mettre en évidence des informations cruciales.

\subsection{Boîte d'Information (\texttt{infobox})}
\begin{infobox}
    \textbf{Astuce :} Ceci est une boîte d'information générée avec \ci{\textbackslash begin\{infobox\}}. Elle utilise les couleurs \texttt{accentPale} pour le fond et \texttt{accent} pour la bordure. Elle est idéale pour les remarques, astuces ou compléments d'information.
\end{infobox}

\subsection{Boîte d'Avertissement (\texttt{warnbox})}
\begin{warnbox}
    \textbf{Attention :} Ceci est une boîte d'avertissement générée avec \ci{\textbackslash begin\{warnbox\}}. Elle utilise des nuances de jaune/brun pour capter l'attention de l'utilisateur sur une action potentiellement destructrice ou une contrainte forte.
\end{warnbox}

\subsection{Note de Marge (\texttt{sidenote})}
\sidenote{Ceci est une note latérale (\texttt{\textbackslash sidenote}). Elle utilise un fond discret \texttt{sidebarBg} et une bordure \texttt{border}. Parfaite pour les métadonnées, les renvois bibliographiques ou des détails techniques très spécifiques.}


% ── SECTION 6 : BLOCS DE CODE ──────────────────────────────────
\section{Blocs de Code}

Le préambule définit trois environnements distincts, basés sur \texttt{tcolorbox} et \texttt{listings}, pour présenter du code source de manière élégante avec une barre latérale colorée.

\subsection{Code Python (\texttt{pyblock})}
\begin{pyblock}
def calculer_trajectoire(vitesse, angle):
    """Calcule la trajectoire en fonction de la vitesse et l'angle."""
    g = 9.81
    distance = (vitesse ** 2) * math.sin(2 * angle) / g
    return distance

# Exemple d'utilisation
print(f"Distance : {calculer_trajectoire(15, 45)}")
\end{pyblock}

\subsection{Code JSON (\texttt{jsonblock})}
\begin{jsonblock}
{
  "project": "Myrtille",
  "version": "1.0",
  "dependencies": {
    "pyside6": "^6.5.0",
    "numpy": ">=1.24"
  },
  "active": true
}
\end{jsonblock}

\subsection{Code Shell (\texttt{shblock})}
\begin{shblock}
#!/bin/bash
# Installation des dépendances
pip install -r requirements.txt
python main.py --config config.json
\end{shblock}


% ── SECTION 7 : TABLEAUX ───────────────────────────────────────
\section{Tableaux}

Les tableaux utilisent les extensions \texttt{booktabs} et \texttt{tabularx} pour un rendu professionnel. L'en-tête du tableau est coloré avec \texttt{tableHdr}.

\vspace{10pt}
\begin{tabularx}{\textwidth}{L{3cm} L{3cm} X}
    \toprule
    \rowcolor{tableHdr}
    \textbf{Paramètre} & \textbf{Type} & \textbf{Description} \\
    \midrule
    \texttt{host} & \badgeStr ~ \badgeReq & Adresse IP ou nom d'hôte du serveur cible. \\
    \texttt{port} & \badgeInt ~ \badgeOpt & Port de connexion (par défaut : 8080). \\
    \texttt{timeout} & \badgeFloat ~ \badgeOpt & Délai d'attente maximum en secondes avant l'abandon. \\
    \texttt{matrix\_data} & \badgeMat ~ \badgeReq & Jeu de données sous forme de matrice bi-dimensionnelle. \\
    \bottomrule
\end{tabularx}